% ===========================================================================
% Section 3: Data
%
% Per Shapiro Robot-mode writing conventions:
%   - State facts; do not argue. Interpretation belongs in Section 5.
%   - Define every concept before using it.
%   - No forward references to undefined constructs.
%   - Plain, precise language.
%
% All numerical claims in this section trace to either:
%   - code/base/*/clean_*.R                (raw network / census totals)
%   - code/pipeline/03a_build_cost_raster.R (cost-surface parameters)
%   - code/pipeline/04_market_access.R     (MA computation)
%   - data/derived/05_panel/departments_wide_panel.parquet (panel stats)
%
% Numbers marked "FROM PANEL" are pulled directly from the estimation
% panel's manifest; those marked "FROM NETWORK" come from
% rails_by_district.parquet or roads_by_district.parquet.
%
% Panel-derived statistics (population means/totals, MA distribution,
% counterfactual component means, 1947 coverage) are AutoFill macros
% from results/scalars.tex. Network-geometry numbers (km totals, segment
% counts) remain inline: they derive from stable raw shapefiles, not
% from the estimation pipeline.
% ===========================================================================

\section{Data}
\label{sec:data}

The empirical analysis combines three data components. Section~\ref{ssec:networks}
describes the transport network geometries and their changes between
1960 and 1986. Section~\ref{ssec:census} describes the district-level
population and economic outcomes. Section~\ref{ssec:ma} describes how
we construct market access from transport cost surfaces.

\subsection{Transport networks}
\label{ssec:networks}

We combine two georeferenced networks: the railway network at its 1979
status, which we map to 1960 and 1986 via each segment's recorded closure
status, and the paved-and-gravel road network at three cross sections
(1954, 1970, 1986).

\paragraph{Railways.}
The rail network comes from the network map in Argentina's 1979
\emph{Política Ferroviaria} report \citep{argentina1981politica},
which we digitised against official IGN geographic
references.\footnote{The report's Figure~V--5 was scanned,
georeferenced, and hand-traced in GIS; the resulting file contains
\railNSegments{} line segments. District-level rail kilometres are
computed by intersecting the segments with the district polygons, which
changes the network total by \railClipLossPct{}~percent.} This is a single
1979 snapshot: the project has no rail map at a second date, and the
periods below are reconstructed from it. Each digitised
segment carries the information needed to reconstruct the network and
to build the instrument. First, a three-way closure status: whether the
segment remained active through 1986, was closed before 1976, or was
closed during the military dictatorship (1976--1983). Two cross sections follow from the field's own
definitions: the network active in 1986, and the network standing before
the dictatorship closures. A third, 1960, requires an assumption,
because the status is undated: that the segments recorded as closed
before 1976 were still open in 1960. Second, its Larkin
Plan status: whether the segment was among the \railNStudied{} segments the
1962 Larkin Plan selected for detailed study --- segments falling
below the Plan's freight-density thresholds, whose closure it then
evaluated on explicit revenue and operating-cost grounds --- or among
the \railNNotStudied{} segments the Plan did not evaluate. For studied segments, the
digitisation also records which of the Plan's three recommendations
(maintain, close, or perform a new study) was issued.
The studied/non-studied distinction is the basis of the first
instrument (Section~\ref{sec:identification}).

Total rail kilometres fall from \railKmSixty{}
in 1960 to \railKmEightySix{} in 1986, a \railPctFall{}\% reduction. Of
the \railKmLost{}~km lost, \railKmDictatorship{}~km
(\railShareDictatorship{}\%) were closed during the dictatorship and
\railKmClosedPreSeventySix{}~km
(\railShareClosedPreSeventySix{}\%) before 1976.
Figure~\ref{fig:networks}a shows the
geography of the rail network and highlights the dictatorship closures
in red. \railNDistrictsNoRail{} of the \nDistricts{} districts have no
rail in any period;
these are concentrated in Patagonia and Tierra del Fuego.

\paragraph{Roads.}
The road network comes from a digitised comparison of
Automóvil~Club~Argentino road maps for 1954, 1970, and 1986
\citep{acamaps}. Each of the 1{,}741 digitised segments is classified
by the set of map years in which it appears, which lets us
reconstruct the paved-and-gravel network at each of the three cross
sections.\footnote{Segments that appear on the 1954 and 1986 maps but
not on the 1970 map are treated as continuously present (a
cartographic omission on the middle map). Segments absent from the
1986 map are excluded from the period totals.}
The network grows from 27{,}513~km in 1954 to 79{,}820~km in 1986, a
190\% increase.
Figure~\ref{fig:networks}b shows the road network and highlights the
segments added between 1954 and 1986 in red.
Table~\ref{tab:network_changes} summarises the network totals by
period.

\begin{figure}[htbp]
\centering
% Raster (PNG) versions of the heavy vector maps are compiled into the
% paper to keep the committed paper.pdf small (vector originals stay in
% results/figures/); see .kiro/tasks.md PDF-bloat entry.
\includegraphics[width=\textwidth]{figure_1_network_changes.png}
\caption{Transport network changes, 1960--1986.}
\label{fig:networks}
\end{figure}

\begin{table}[htbp]
\centering
\caption{Transport Network Changes, 1960--1986}
\label{tab:network_changes}
\input{../results/tables/table_1_network_changes.tex}
\end{table}

Three districts, all in Patagonia, lack any segment in the road
shapefile and contribute missing values in the district-level road
variables. 11 districts lost rail without gaining roads and 141 districts
both lost rail and gained roads. The cross-district correlation between
$\Delta$(rail km) and $\Delta$(road km) is $-0.25$.

\paragraph{Hypothetical road networks.}
We also construct four hypothetical road networks that connect
Argentina's largest cities under stylised routing rules. The node set
is the 54~cities that were provincial capitals or had population of
at least 15{,}000 in 1960.\footnote{Capital Federal is not among the
nodes, although it does enter every district's market access as a
destination. The asymmetry is deliberate rather than an omission: the
market-access weights ask which destinations a district could reach,
and access to Buenos~Aires counts for all of them, whereas the
hypothetical network asks which intercity links a road programme could
have built, and Buenos~Aires is the terminus those links run to rather
than one of the cities to be connected.} The first network links every pair of
cities with a straight line (1{,}431 segments); the second is the
minimum spanning tree of those straight-line distances (53 edges).
The third replaces straight lines with least-cost paths routed over a
construction-cost surface that, following \citet{faber2014},
penalises slope, wetlands, urban land, and water crossings; the
fourth is the minimum spanning tree of those least-cost-path weights.
These networks serve as instruments for the actual road network in
Section~\ref{sec:identification}; the least-cost-path spanning tree
is the main choice, and the other three appear as robustness
alternatives.

\subsection{Census data and outcomes}
\label{ssec:census}

\paragraph{Geographic unit.}
We work at the district (\emph{departamento}) level using IPUMS
International's time-invariant district geography for Argentina
\citep{ipums2020}, which partitions the mainland into 312
districts, including Capital Federal and the two Tierra~del~Fuego
jurisdictions. The same polygons apply across the five IPUMS sample
years (1970, 1980, 1991, 2001, 2010), the 1947 and 1960 census
publications, and the agricultural and industrial censuses, so every
outcome is measured on identical boundaries throughout.

\paragraph{Population.}
Population in 1960 comes from the digitised \emph{Censo Nacional de
Población 1960} tables \citep{censo1960pop}, reconciled to the IPUMS district geography
through a crosswalk that handles pre-1960 splits and name changes. The
mean 1960 district population is \popMeanBase{} with a national total of
\popTotalBase{}~million. Population in 1991 comes from the IPUMS microdata
aggregated to the same district geography; the mean 1991 district population is
\popMeanEnd{} with a national total of \popTotalEnd{}~million. Log population change
$\Delta\log(\mathrm{pop}_{91/60})$ has mean $\chgLogPopMean{}$ and standard deviation
$\chgLogPopSD{}$. Capital Federal and the two Tierra~del~Fuego districts are
covered by both census years; Capital Federal contributes to the
market-access population weights defined in Section~\ref{ssec:ma} but
is excluded from the estimation sample as an observation
(Section~\ref{sec:results}).

\paragraph{Urban population.}
Urban population shares at the district level come from the 1960
census (classification per INDEC) and from the IPUMS 1991 microdata.
Urban population is not coded consistently in the 1970 and 2010 IPUMS
microdata for Argentina; those years are recorded as missing in the
panel and are not used for urban-outcome regressions. 1960 and 1991 give
us the post-reform window with consistent urban/rural classification.

\paragraph{Pre-reform placebo.}
We supplement the baseline with the 1947 national census
\citep{censo1947}, which covers
\placeboCoverage{} of the \nDistricts{} districts. A log population change
$\Delta\log(\mathrm{pop}_{60/47})$ for this subsample serves as a placebo
outcome: it measures population growth before any of the 1960--1986
network changes, and should be uncorrelated with the instruments if the
instruments isolate post-1960 network-induced variation.

\paragraph{Industrial and agricultural outcomes.}
Sectoral outcomes come from four additional censuses. The 1954 and
1985 industrial censuses (\emph{Censo Industrial} and \emph{Censo
Económico}; \citealt{censo1954industrial, indec1985economico}) provide
district-level establishment counts, total wages paid (the wage
mass), and production values, each measured on definitions comparable
across the two census years. Employment is recorded on incompatible
definitions in the two years (workers in 1954, total persons engaged
in 1985) and is not used. The 1960
and 1988 agricultural censuses
\citep{censo1960agro, indec1988agro} provide the number of farm
operations and total agricultural area. All monetary
variables are nominal pesos of their census year; real comparisons
across years require a deflator that we defer to robustness.

\paragraph{Geographic controls.}
District-level controls come from external sources. Area in~km$^2$ is
computed from the IPUMS polygon. Mean elevation comes from GTOPO30
\citep{usgs1996gtopo30}. Terrain ruggedness uses the Nunn--Puga TRI
raster \citep{nunn2012ruggedness} derived from GTOPO30. Wheat potential
yield comes from FAO--GAEZ~3.0 \citep{faoiiasa2012gaez}. Caloric
potential (pre-~and post-1500) comes from \citet{galorozak2016}.
Distance to Buenos Aires is the great-circle distance from the centroid
of the Gran~Buenos~Aires polygon in the IGN settlements shapefile
\citep{ignsettlements}.
All continuous controls are standardised (zero mean, unit variance) for
use in the regressions. Appendix Table~\ref{tab:descriptives} reports
descriptive statistics for every variable entering the estimation.

\subsection{Market access}
\label{ssec:ma}

\paragraph{Definition.}
Following \citet{donaldsonhornbeck2016},
we define district $i$'s market access in period $t$ for sector
$s$ under elasticity $\theta$ as:
\begin{equation}
  \mathrm{MA}_{i,t,s,\theta} = \sum_{j \neq i}
     \frac{\mathrm{pop}_{j,1960}}{(\tau_{ij,t,s})^{\theta}},
  \label{eq:ma}
\end{equation}
where $\tau_{ij,t,s}$ is the origin--destination transport cost between
the centroids of districts~$i$ and~$j$ under network configuration~$t$
and sector~$s$. Population weights are held fixed at their 1960 values
across all $t$ so that variation in~$\mathrm{MA}_{i,t,s,\theta}$ comes
entirely from variation in transport costs rather than from endogenous
population responses to the network changes. All 312 districts
contribute to the weights, including Capital Federal
(2.97~million people in 1960): access to Buenos~Aires enters every
district's market access even though Capital Federal itself is not an
observation in the estimation sample (Section~\ref{sec:results}).

\paragraph{Cost surface.}
Transport costs $\tau_{ij,t,s}$ are the least-cost-path distances on
a $2{,}399 \times 3{,}090$ raster covering mainland Argentina in the
ESRI:54034 equal-area projection (approximately 1.2~km per cell). The
per-pixel cost $c_i(t,s)$ is:
\begin{equation*}
c_i(t, s) =
\begin{cases}
\text{cost}_{\mathrm{nav}}(s)                      & \text{if pixel $i$ is navigable water,} \\
\text{cost}_{\mathrm{mininfra}}(s)                 & \text{if pixel $i$ has both rail and road,} \\
\text{cost}_{\mathrm{rail}}(s)                     & \text{if pixel $i$ has rail only,} \\
\text{cost}_{\mathrm{road}}(s)                     & \text{if pixel $i$ has road only,} \\
\text{cost}_{\mathrm{land}} \cdot \mathrm{HMI}_i   & \text{if pixel $i$ is off-network land in Argentina,} \\
\text{NA}                                          & \text{otherwise (hard barrier).} \\
\end{cases}
\end{equation*}
Network layers (rail, road, and navigation) are rasterised from the
respective shapefiles after a 1~km buffer to guarantee at least one
full raster cell per linear segment; the navigation layer comprises
the watercourse segments flagged navigable in the IGN hydrography
shapefile \citep{ignhydro}. The Strait of Magellan is buffered
by 5~km so that at least one raster cell per shore becomes navigable,
which lets Dijkstra cross the strait and keeps Tierra~del~Fuego
connected to the mainland. Off-network pixels outside Argentina are
set to NA, which acts as a hard barrier for the least-cost path
algorithm except where a navigable-water pixel overrides the country
mask (e.g.\ the Paraná river border).
Figure~\ref{fig:navigation} maps the navigation layer: the navigable
set is the inland Paraná--Paraguay--Uruguay--Río de la Plata system
plus a few Patagonian rivers, together with the Magellan crossing.
The cost surface contains no open-ocean coastal shipping, so
Atlantic ports enter the network through land or river legs only.

\paragraph{Cost parameters.}
Sector-specific cost parameters come from \citet{baumgartnerpalazzo1969}
Table~II. Three sectors $s \in \{0, 1, 2\}$ correspond to three cargo
density scenarios (medium, high, low); Table~\ref{tab:cost_params}
summarises the values.

\begin{table}[htbp]
\centering
\small
\caption{Transport Costs by Mode and Cargo Density
         (1960 pesos per ton-km)}
\label{tab:cost_params}
\begin{tabular}{lccc}
\toprule
              & Low density  & Medium density & High density \\
              & (100 t/day)  & (500 t/day)    & (1{,}000 t/day) \\
\midrule
Road          & 8.266        & 1.777          & 1.000 \\
Railroad      & 13.490       & 1.874          & 0.465 \\
Navigation    & 0.621        & 0.621          & 0.621 \\
\bottomrule
\end{tabular}

\vspace{0.5em}
{\footnotesize
\begin{minipage}{0.85\textwidth}
\emph{Notes}: Ground-transport costs from
\citet{baumgartnerpalazzo1969}. Navigation costs from
\citet{larkin1962}. Medium-density costs are used in the main
analysis (sector $s = 0$); the headline population result is
re-estimated under the high- and low-density schedules (sectors
$s = 1$ and $s = 2$) in Table~\ref{tab:density_schedules}.
\end{minipage}
}
\end{table}

In particular, rail is cheaper than road at high cargo density
($\text{cost}_\mathrm{rail}(s{=}1) = 0.465 <
\text{cost}_\mathrm{road}(s{=}1) = 1.000$) but more expensive at low
density ($\text{cost}_\mathrm{rail}(s{=}2) = 13.490 >
\text{cost}_\mathrm{road}(s{=}2) = 8.266$), which is the
scale-economies channel that motivates sector-specific analysis. At
medium density, rail and road costs are nearly identical (1.874 vs.\
1.777), so mode choice matters less for the overall market-access
index. We therefore use the medium-density scenario as the main
specification --- it is the general-purpose schedule and does not
build a mode preference into the headline treatment --- and report
the population estimates under all three schedules alongside the
main results (Table~\ref{tab:density_schedules}). The
off-network cost $\text{cost}_{\mathrm{land}}$ is
$17.9 \times \text{cost}_\mathrm{road}(s{=}2) / \overline{\mathrm{HMI}}$,
where $\overline{\mathrm{HMI}} = 0.2018$ is the national mean of the
Özak~\citeyearpar{ozak2018} Human Mobility Index. The factor~17.9 is
the ratio of motor-vehicle speed to walking-equivalent speed, calibrated
on a Buenos~Aires--Rosario comparison. This gives
$\text{cost}_{\mathrm{land}} \approx 733.2$; the per-pixel off-network
cost scales with the local HMI value and is one to two orders of
magnitude above on-network cost.

\paragraph{Least-cost path and elasticity.}
We compute $\tau_{ij,t,s}$ by running Dijkstra's algorithm over the
cost surface between the 312 district centroids, separately for each
network configuration and sector.\footnote{Paths move across an
eight-connected grid; the computation uses the \texttt{gdistance}
R~package. The replication package documents runtimes.}
$\tau_{ii,t,s}=0$ by construction and $\tau$ is symmetric.

The elasticity $\theta$ controls how sharply distant destinations are
discounted in the MA sum. It is taken from the trade literature, not
estimated in this paper. Structural
estimates span a wide range: \citet{eatonkortum2002} report a
preferred value of 8.28; \citet{simonovskawaugh2014} revise the
Eaton--Kortum approach with simulated-moments estimation and obtain
values near 4, with a credible range of roughly 4 to 8;
\citet{caliendoparro2015} estimate 8.11 for agriculture; and
\citet{donaldson2018} estimates 3.8, the only estimate from
railroad-borne trade (colonial India). We report
two values within this range, each taken from a single study:
$\theta_{\mathrm{low}} = \thetaLow{}$, the \citet{simonovskawaugh2014}
benchmark estimate (their exactly identified simulated-moments estimate
on 2004 price and trade-flow data for 123 countries, standard error
\thetaLowSE{}), and $\theta_{\mathrm{high}} = \thetaHigh{}$, the
\citet{caliendoparro2015} agricultural trade elasticity. Main results use
$\theta_{\mathrm{low}}$; $\theta_{\mathrm{high}}$
appears in the sensitivity tables, and
Section~\ref{sec:results_robustness} reports the estimates over the
full sweep $\theta \in [1, 12]$ (Table~\ref{tab:theta_sweep});
Section~\ref{sec:discussion} discusses what the sweep implies for
comparisons with the literature.
\placeholder{whether a trade-elasticity exponent is the right object
for our raw generalized cost $\tau$ is memo Decision~A, unresolved}
% THETA PROVENANCE. 8.11 traces to Caliendo-Parro's agricultural trade
% elasticity (via D&H 2016 fn. 55). The pipeline inherited theta_low =
% 4.55 with no traceable source: a provenance search (2026-07-16) found
% no occurrence in Simonovska-Waugh 2014, Caliendo-Parro 2015, D&H 2016
% fn. 55, Fajgelbaum-Redding 2022, or the old draft; Cote (2026-07-24)
% did not recall one, and 4.55 was carried as a declared midpoint.
% RESOLVED 2026-09 (Cote email, Diego approval): theta_low is now the
% Simonovska-Waugh published benchmark 4.14 (SE 0.09; their Table 5 and
% abstract; the NBER revision said 4.12 and D&H fn. 55 quote the
% overidentified 4.10). Every output regenerated; see config.R section
% 8. Memo Decision A (whether a trade-elasticity exponent belongs on a
% raw, non-normalized tau at all) is NOT resolved by citing the source
% and stays in the placeholder above.

\paragraph{Change in log MA.}
The main treatment variable is
$\Delta\log\mathrm{MA}_i \equiv \log\mathrm{MA}_{i,1986,s=0,\theta_{\mathrm{low}}}
  - \log\mathrm{MA}_{i,1960,s=0,\theta_{\mathrm{low}}}$.
Across the 312~districts,
$\Delta\log\mathrm{MA}_i$ has mean $\maMean{}$, standard deviation $\maSD{}$,
minimum $-\maMin{}$ and maximum $\maMax{}$; \maSharePos{}\% of districts see a positive
change. Figure~\ref{fig:ma_choropleth} maps $\Delta\log\mathrm{MA}_i$
across districts.

\begin{figure}[htbp]
\centering
\includegraphics[width=0.85\textwidth]{figure_2_ma_change_choropleth.png}
\caption{Change in log market access, 1960--1986.}
\label{fig:ma_choropleth}
\end{figure}

\paragraph{Counterfactual decomposition.}
We additionally construct two counterfactual cost surfaces that freeze
one network mode at its baseline while letting the other change:
\begin{itemize}\setlength\itemsep{0pt}
  \item \emph{only-rail}: 1986 rails combined with 1954 roads.
  $\Delta\log\mathrm{MA}^{\text{only-rail}}_i = \log\mathrm{MA}(\text{only-rail})_i
    - \log\mathrm{MA}_{i,1960}$.
  \item \emph{only-road}: 1960 rails combined with 1986 roads.
  $\Delta\log\mathrm{MA}^{\text{only-road}}_i$ defined analogously.
\end{itemize}
At sector~0 and $\theta_{\mathrm{low}}$, $\Delta\log\mathrm{MA}^{\text{only-rail}}$
has mean $-\cfRailMean{}$ (rail contraction is a small market-access loss) and
$\Delta\log\mathrm{MA}^{\text{only-road}}$ has mean $\cfRoadMean{}$ (road
expansion is a large gain). The sum of the two isolated components
($\cfSumComponents{}$) is slightly below the total
$\Delta\log\mathrm{MA} = \maMean{}$, reflecting a positive complementarity
between the two network shocks. Figure~\ref{fig:ma_counterfactual}
compares the three maps side by side.

\begin{figure}[htbp]
\centering
\includegraphics[width=\textwidth]{figure_c13_ma_counterfactual_trio.png}
\caption{Market access changes under three scenarios: full, only-rail
counterfactual, only-road counterfactual.}
\label{fig:ma_counterfactual}
\end{figure}

\paragraph{Limitations.}
The cost-surface construction imposes three modelling choices that the
Section~\ref{sec:identification} discussion of identification returns to.
First, mode switches are free: Dijkstra can step from a rail pixel to a
road pixel at any grid boundary, whereas in reality such switches
occur at stations or terminals and incur transshipment costs.
Section~\ref{sec:results_robustness} bounds this assumption from the
opposite extreme (one mode per trip, i.e.\ infinite transshipment
cost). Second,
the algorithm selects a single cheapest path, not a probability
distribution over routes; in equilibrium firms use multiple parallel
routes. Assigning each cell the cheapest mode deterministically also
treats modes as perfect substitutes along the path;
\citet{fuchswong2026} estimate an elasticity of modal substitution
near~1.1 between rail and truck in the contemporary United States,
far from perfect substitution.
Third, cost parameters are constant across space and time; we
do not model congestion or regional cost heterogeneity. Each of these
simplifications follows the approach in
\citet{donaldsonhornbeck2016} and is common in the market-access
literature.
